\documentclass[10pt,a4paper]{amsart}
\usepackage[margin=19mm]{geometry}
\usepackage{amsmath,amssymb,mathtools}
\usepackage[scr=rsfs,cal=euler]{mathalfa}
\usepackage{xfrac}
\usepackage[unicode,hidelinks,bookmarks=false]{hyperref}
\newcommand{\ie}{i.e.\ }
\newcommand{\cf}{cf.\ }
\newcommand{\resp}{respectively\ }
\newcommand{\Subsection}{\subsection}
\providecommand{\nobreakdash}{\nobreak\hbox{-}}
\setlength{\emergencystretch}{2em}
\multlinegap0pt
\usepackage{siunitx}
\usepackage{siunitx}
\usepackage{siunitx}
\usepackage{siunitx}
\usepackage{amssymb}
\usepackage{mathtools}
\usepackage{xcolor}
\usepackage{tikz}
\usepackage{siunitx}
\newtheorem{lemma}{Lemma}
\newcommand{\abs}[1]{\left\lvert #1 \right\rvert}
\newcommand{\bm}[1]{\boldsymbol{#1}}
\newcommand{\mc}[1]{\mathcal{#1}}
\newcommand{\norm}[1]{\left\lVert #1 \right\rVert}
\newcommand{\paren}[1]{\left(#1\right)}
\title{A Cartesian Grid Method for Advection-Diffusion Equations with Robin Boundary Conditions on Moving Domains}
\author{Han Zhou, Yoichiro Mori, Lingxing Yao}
\date{Source: arXiv:2607.26341v1 (CC BY 4.0)}
\begin{document}
\maketitle

\noindent\emph{Source and licence.} A Cartesian Grid Method for Advection-Diffusion Equations with Robin Boundary Conditions on Moving Domains. Version arXiv:2607.26341v1 (CC BY 4.0), distributed by arXiv under the Creative Commons Attribution 4.0 International licence, http://creativecommons.org/licenses/by/4.0/, as recorded in the arXiv licence metadata for this exact version and shown on the abstract page https://arxiv.org/abs/2607.26341v1. Full licence notice: \texttt{licenses/arxiv-2607.26341-CC-BY-4.0.txt}.\par\smallskip

\noindent\textit{Editorial scope.} Counted unit: the article's lemma lem:KM-mat-o2 (source0.tex lines 893--907). The article's complete original proof of it is delivered with it (source0.tex lines 909--1047, one \texttt{proof} environment of 139 lines). No mathematics is written, repaired or re-derived: the statement and the proof are the article's own lines, in the article's own notation, and no retained line is substituted or re-flowed.  No further source-local context is needed: every cross-reference of every retained line resolves inside the two delivered spans.  Nothing was omitted from any retained span, so the package carries no omission record.  No retained line contains a citation, so the wrapper carries no bibliography and no bibliography entry.  No retained line contains a figure environment, an image-inclusion command or drawing code, so no image is delivered and the package declares no asset dependency.  Editorial formatting only, in addition: an AMS \texttt{amsart} preamble that restates the article's own declarations and macros used by the retained lines, namely the environments \texttt{lemma} on separate counters, together with the standard packages those lines need.  The retained environments print in this document as Lemma 1 in order of appearance, whereas the article prints the same items with the numbering of its own full document.  The complete native source of the article as distributed in the arXiv e-print for this exact version is bundled unchanged as \texttt{sources/206261/source0.tex}.
\begin{lemma}\label{lem:KM-mat-o2}
Assume $h<2/\abs{v}$ when $v\ne0$ and $1-\lambda\tau>0$.
For sufficiently small $h$ and $\tau$, the following
statements hold for $n = 1,2,\cdots, N_t-1$:
\begin{enumerate}
    \item The matrix $K_I^n$ is an M-matrix and satisfies
\begin{equation}
    \norm{(K_I^n)^{-1}}_\infty \le 1+C\tau,
\end{equation}
    \item The bound $\norm{T^n}_\infty \le \frac{5}{3}$ holds. If $\tau / h^2 \ge 1/2$, the matrix $P^n$ has nonnegative entries and satisfies
\begin{equation}
    \norm{P^n}_\infty \le 1+C\tau.
\end{equation}
\end{enumerate}
\end{lemma}

\begin{proof}
To prove (1), we consider the entries of the matrix $K_w^n$ since $K_I^n$ is the top-left $j(n)\times j(n)$ block of $K_w^n$.
The matrix $\frac{1}{Q^n}E^n(D^n)^T$ is nonzero only for the block $j(n):j(n)+1, j(n)-1:j(n)+1$.
So, for $i \neq j(n)$, assuming $h < 2 / \abs{v}$, we have
\begin{align}
    (K_w^n)_{i,i-1} &= -\paren{\frac{\tau}{h^2} + \frac{v \tau}{2 h} } e^{-\beta h} < 0, \\
    (K_w^n)_{i,i+1} &= -\paren{\frac{\tau}{h^2} - \frac{v \tau}{2 h} } e^{\beta h} < 0, \\
    (K_w^n)_{ii} &= 1 + \frac{2\tau}{h^2} > 0.
\end{align}
For sufficiently small $h$,
\begin{equation}
\begin{aligned}
    (K_w^n)_{i,i-1} + (K_w^n)_{i,i} + (K_w^n)_{i,i+1}
    &= 1 + \frac{2\tau}{h^2} \paren{1 - \cosh(\beta h)} + \frac{v \tau}{h} \sinh(\beta h)  \ge 1 - C \tau.
\end{aligned}
\end{equation}
We consider the $j(n)$-th row of $K_I^n$,
\begin{equation}
\begin{aligned}
    (K_w^n)_{j(n),j(n)-1} &= e^{-\beta h}\frac{\tau}{h^2}\paren{ -1 - \frac{vh}{2} + \frac{\paren{1-\frac{vh}{2}} \paren{\sigma^n-\frac{1}{2}} }{\paren{\sigma^n+\frac{1}{2}+\sigma^n\alpha^nh} } },
\end{aligned}
\end{equation}
\begin{equation}
\begin{aligned}
    (K_w^n)_{j(n),j(n)} &=  1 + \frac{\tau}{h^2}\paren{2 + \frac{\paren{1-\frac{vh}{2}} \paren{-2\sigma^n+\paren{1-\sigma^n}\alpha^n\sigma^nh} }{\paren{\sigma^n+\frac{1}{2}+\sigma^n\alpha^nh} }}.
\end{aligned}
\end{equation}
Then, we have
\begin{equation}
\begin{aligned}
    (K_w^n)_{j(n),j(n)-1} + (K_w^n)_{j(n),j(n)}  \ge 1 + \frac{2\tau}{h} \paren{  \frac{\beta + \alpha^n\sigma^n(2-\sigma^n)}{2\sigma^n+1} }- C\tau\ge 1 - C\tau,
\end{aligned}
\end{equation}
where we have used
\begin{equation}
    \beta + \alpha^n\sigma^n(2-\sigma^n) \ge \beta + \alpha_0\sigma^n(2-\sigma^n) \ge 1 +\abs{\alpha_0} - \abs{\alpha_0}  \ge 1.
\end{equation}
Thus, for sufficiently small $h$, we have
\begin{equation}\label{eqn:kw-rowsum}
    K_I^n\bm 1 \ge \paren{1 - C\tau}\bm 1,
\end{equation}
where $C>0$ depends on $v$ and $\alpha$.
Since the diagonal elements are positive and the off-diagonal elements are non-positive for sufficiently small $h$ and $\tau$, \eqref{eqn:kw-rowsum} implies that $K_I^n$ is an M-matrix and satisfies
\begin{equation}
    \norm{(K_I^n)^{-1}}_\infty \le 1 + C\tau.
\end{equation}





We now prove (2).
If $j(n)-1 \neq j(n-1)$, the entries of the matrix $T^n=R^n M_w^{n-1} (R^{n-1})^T$ are nonzero only along the diagonal and are all equal to $1$, yielding $\norm{T^n}_\infty=1$.
In the case $j(n)-1 = j(n-1)$, we have
\begin{align}
    \paren{M_w^{n-1}}_{j(n), j(n)-2} &= e^{-2\beta h}\frac{\frac{1}{2}-\sigma^{n-1}}{\sigma^{n-1}+\frac{1}{2}+\sigma^{n-1}\alpha^{n-1}h},\\
    \paren{M_w^{n-1}}_{j(n), j(n)-1} &= e^{-\beta h}\frac{2\sigma^{n-1} - \paren{1-\sigma^{n-1}} \alpha^{n-1} h }{ \frac{1}{2} + \sigma^{n-1} + \sigma^{n-1} \alpha^{n-1} h }.
\end{align}
For $\sigma^{n-1} \in[0,1/2]$ and sufficiently small $h$, if $\alpha^{n-1}\ge 0$, then $\paren{M_w^{n-1}}_{j(n), j(n)-2} \ge 0$, and we have
\begin{equation}\label{eqn:Mwrowjn}
\begin{aligned}
    &\abs{\paren{M_w^{n-1}}_{j(n), j(n)-2}} + \abs{\paren{M_w^{n-1}}_{j(n), j(n)-1}} \\
    &= \frac{ e^{-2\beta h}\paren{\frac{1}{2}-\sigma^{n-1}} + e^{-\beta h}\paren{ 2\sigma^{n-1} - \paren{1-\sigma^{n-1}} \alpha^{n-1} h }}{ \frac{1}{2} + \sigma^{n-1} + \sigma^{n-1} \alpha^{n-1} h } \\
    &\le 1 - \frac{\paren{\beta+\alpha^n}h + \mc O(h^2) }{ \frac{1}{2} + \sigma^{n-1} + \sigma^{n-1} \alpha^{n-1} h } \le 1.
\end{aligned}
\end{equation}
If $\alpha^{n-1}<0$, then it is possible that $\paren{M_w^{n-1}}_{j(n), j(n)-2} < 0$, in which case we have
\begin{equation}
\begin{aligned}
    &\abs{\paren{M_w^{n-1}}_{j(n), j(n)-2}} + \abs{\paren{M_w^{n-1}}_{j(n), j(n)-1}} \\
    &= \frac{ e^{-2\beta h}\paren{\frac{1}{2}-\sigma^{n-1}} - e^{-\beta h}\paren{ 2\sigma^{n-1} - \paren{1-\sigma^{n-1}} \alpha^{n-1} h }}{ \frac{1}{2} + \sigma^{n-1} + \sigma^{n-1} \alpha^{n-1} h } \\
     &\le \frac{ e^{-2\beta h}\paren{\frac{1}{2}-\sigma^{n-1}} + e^{-\beta h} \alpha^{n-1} h }{ \frac{1}{2} + \sigma^{n-1} + \sigma^{n-1} \alpha^{n-1} h } \\
     &= \frac{ \frac{1}{2}-\sigma^{n-1} + \paren{\alpha^{n-1} - (1-2\sigma^{n-1})\beta}h + \mc O(h^2) }{ \frac{1}{2} + \sigma^{n-1} + \sigma^{n-1} \alpha^{n-1} h }  \le 1.
\end{aligned}
\end{equation}
This implies that for $\sigma^{n-1} \in[0,1/2]$, we also have $\norm{T^n}_\infty \le 1$  and
\begin{equation}
    \norm{P^n}_\infty \le \norm{ R^{n+1} M_w^{n}  (R^{n})^T}_\infty \norm{(K_I^n)^{-1} }_\infty \le 1 + C\tau.
\end{equation}
If $\sigma^{n-1} \in (1/2,1)$, the only row whose absolute row sum can exceed
one is again the moving-interface row. For sufficiently small $h$, the second
entry in that row is positive, and therefore
\begin{equation}
\begin{aligned}
    \sum_k\abs{(T^n)_{j(n),k}}
    &\le
    \frac{ e^{-2\beta h}\paren{\sigma^{n-1}-\frac12}
    + e^{-\beta h}\abs{2\sigma^{n-1}-(1-\sigma^{n-1})\alpha^{n-1}h}}
    {\frac12+\sigma^{n-1}+\sigma^{n-1}\alpha^{n-1}h}  \\
    &=
    \frac{3\sigma^{n-1}-\frac12}{\sigma^{n-1}+\frac12}+\mc O(h)
    \le \frac53 .
\end{aligned}
\end{equation}
The last inequality is uniform for sufficiently small $h$. Hence $\norm{T^n}_\infty\le5/3$.


Now, we consider the matrix $P^n=R^{n+1} M_w^{n} (R^{n})^T(K_I^n)^{-1}$.
If $\sigma^{n}\in[0,1/2]$, the non-negativity and norm estimate follows immediately from the preceding row-sum bound and the M-matrix inverse bound.

For $\sigma^n \in (1/2,1)$, let
$\bm z=R^{n+1}M_w^n(R^n)^T\bm w$ and
$\bm w=(R^nK_w^n(R^n)^T)^{-1}\bm y$, where
\begin{equation}
    R^n K_w^n (R^n)^T \bm w = \bm y.
\end{equation}
Since $(K_I^n)^{-1}$ is non-negative and $R^{n+1} M_w^{n} (R^{n})^T$ is nonzero only along the diagonal and equals $1$ for rows from $1$ to $j(n)$, we only need to check the $(j(n)+1)$-th row of $P^n$. Here, $c_{j(n)+1}$ is a ghost-node value on the box $\mc B$, as in the interface formulation.
Since $c_j = w_j$ for $j \le j(n)$, the ghost value $c_{j(n)+1}$ satisfies
\begin{align}
    &e^{-\beta h}\paren{\sigma^n-\frac{1}{2}} w_{j(n)-1} + \paren{-2\sigma^n + \paren{1-\sigma^n}\alpha^n h}w_{j(n)} +e^{\beta h} \paren{\sigma^n + \frac{1}{2} + \sigma^n \alpha^n h} c_{j(n)+1} = 0,\\
    &- e^{-\beta h}\paren{\frac{\tau}{h^2}+\frac{v\tau}{2h}}w_{j(n)-1} + \paren{1+2\frac{\tau}{h^2}} w_{j(n)}  - e^{\beta h} \paren{\frac{\tau}{h^2} - \frac{v\tau}{2h}} c_{j(n)+1} = y_{j(n)}.
\end{align}
By eliminating $w_{j(n)-1}$, we obtain
\begin{align}
c_{j(n)+1} &= e^{-\beta h} \frac{\eta_1 y_{j(n)} + \eta_2 w_{j(n)}}{\eta_3}, \quad \eta_1 = \sigma^n-\frac{1}{2}, \\
    \eta_2 &= \frac{1}{2}-\sigma^n + \frac{\tau}{h^2} \paren{ 1-(1-\sigma^n)\alpha^nh+\frac{vh}{2}\Bigl(2\sigma^n-(1-\sigma^n)\alpha^nh\Bigr)},  \\
    \eta_3 &= \frac{\tau}{h^2}\paren{ 1+\sigma^n\alpha^nh + \frac{vh}{2} \sigma^n \paren{2+\alpha^nh}}.
\end{align}
For $\frac{\tau}{h^2}>\frac{1}{2}$ and sufficiently small $h$, since $\sigma^n\in (\frac{1}{2},1)$, we have
\begin{equation}
    \eta_1, \eta_3 > 0,\quad \eta_2 \ge \frac{1}{2}-\sigma^n + \frac{1}{2}>0.
\end{equation}
This implies the $(j(n)+1)$-th row of $P^n$ is also non-negative and
\begin{equation}
    \abs{c_{j(n)+1}} \le e^{-\beta h}\frac{\eta_1+\eta_2}{\eta_3}\max\{\abs{y_{j(n)}}, \abs{w_{j(n)}}\}.
\end{equation}
Since $e^{-\beta h} = 1 - \beta h + \mc O(h^2)$, by a direct calculation, we have
\begin{equation}
    e^{-\beta h} \frac{\eta_1+\eta_2}{\eta_3} = 1 - (\beta+\alpha^n)h + \mc O(h^2).
\end{equation}
Using the condition $\beta + \alpha^n \ge 1$, we have $e^{-\beta h} \frac{\eta_1+\eta_2}{\eta_3} \le 1 - h + \mc O(h^2) \le 1$ for sufficiently small $h$. This implies
\begin{equation}
    \abs{c_{j(n)+1}} \le \paren{1+C\tau}\norm{\bm y}_\infty,
\end{equation}
and so
\begin{equation}
    \norm{P^n}_\infty \le 1+C\tau.
\end{equation}
\end{proof}

\end{document}
